% ====================================================================
% Type-k Neutrosophic Sets: Recursive Triadic Structures and Their
% Empirical Realization in Large Language Models
%
% Authors: Florentin Smarandache (1), Maikel Leyva-Vázquez (2)
% Target:  Neutrosophic Sets and Systems
% Draft:   v0.1 — 2026-05-04
% ====================================================================

\documentclass[10pt,a4paper]{article}
\usepackage{amsmath,amssymb,amsthm}
\usepackage{booktabs}
\usepackage{hyperref}
\usepackage{graphicx}
\usepackage{array}
\usepackage{xcolor}
\usepackage[margin=2.5cm]{geometry}

% Theorem environments
\newtheorem{definition}{Definition}
\newtheorem{theorem}{Theorem}
\newtheorem{corollary}{Corollary}
\newtheorem{proposition}{Proposition}
\newtheorem{example}{Example}
\newtheorem{remark}{Remark}

\title{\textbf{Type-$k$ Neutrosophic Sets: Recursive Triadic Structures\\
and Their Empirical Realization in Large Language Models}}

\author{
  Florentin Smarandache$^{1}$ \and Maikel Leyva-V\'{a}zquez$^{2}$\\[6pt]
  \small $^{1}$Department of Mathematics, University of New Mexico,
  Gallup, NM 87301, USA. \texttt{smarand@unm.edu}\\
  \small $^{2}$Facultad de Ciencias Matem\'{a}ticas,
  Universidad de Guayaquil, Ecuador.
  \texttt{maikel.leyvavazquez@ug.edu.ec}
}

\date{May 2026}

\begin{document}
\maketitle

% ====================================================================
\begin{abstract}
We introduce \emph{Type-$k$ Neutrosophic Sets}, a recursive extension
of neutrosophic logic in which each epistemic component $(T, I, F)$
is itself characterised by a full neutrosophic triplet, yielding
$3^k$ scalars per element at recursion depth $k$.
We prove a strict expressive hierarchy
$\text{Type-1} \subsetneq \text{Type-2} \subsetneq \cdots \subsetneq \text{Type-}k$
and identify a canonical embedding $\varphi$ that maps any
Type-$(k-1)$ element into its Type-$k$ representation.
We propose a recursive aggregation operator that reduces to the
standard Single-Valued Neutrosophic Weighted Average (SVNWA)
at $k=1$. Empirically, we demonstrate that large language models
(LLMs) implicitly compute Type-2 values when evaluation protocols
permit components outside $[0,1]$: across 2,419 evaluations of
six foundation models under five epistemic protocols,
24.7\% of all responses require Type-2 representation.
Under the overset protocol S4-O.A, this fraction rises to
97.3\%, and is consistent across all six vendors (19--30\%).
The restriction to Type-1 is a \emph{protocol artefact}, not an
epistemic limit. Type-$k$ Neutrosophic Sets provide the first formal
framework that makes these LLM behaviours mathematically legitimate.
\end{abstract}

\textbf{Keywords:} neutrosophic sets; type-2 neutrosophic;
recursive neutrosophic; large language models; epistemic uncertainty;
epistemic indeterminacy; foundation models.

% ====================================================================
\section{Introduction}

Smarandache's neutrosophic logic \cite{smarandache1998neutrosophy}
represents propositions as triplets $(T, I, F)$ where $T$ is the
degree of truth, $I$ the degree of indeterminacy, and $F$ the degree
of falsity, with $T, I, F \in [0,1]$ independently (no sum
constraint). This \emph{Type-1} representation has found application
in multi-criteria decision making \cite{wang2010svn}, image
processing, and, more recently, in the epistemic evaluation of large
language models (LLMs) \cite{mason2026tensors,
leyvavazquez2026plithogenic}.

A fundamental question arises naturally: if $T$ is the degree to
which a proposition is true, to what degree is $T$ \emph{itself}
true, indeterminate, or false? In type-2 fuzzy logic
\cite{mendel2002type2}, membership degrees are themselves fuzzy
sets, enabling a second layer of uncertainty. The neutrosophic
analogue has not been formally defined.

In this paper we fill that gap. We introduce \emph{Type-$k$
Neutrosophic Sets}, where each component of a Type-$(k-1)$
representation is replaced by a full neutrosophic triplet.
The resulting structure has $3^k$ scalars per element
at depth $k$ and subsumes Type-1 as the base case.

Our motivation is both theoretical and empirical. On the theoretical
side, the construction is the most natural extension of neutrosophic
logic to nested uncertainty. On the empirical side, we show that
state-of-the-art LLMs \emph{already compute Type-2 values} when the
evaluation protocol opens the range beyond $[0,1]$: negative $T$
values (``anti-truth''), $I > 1$ (``hyper-indeterminacy''), and
$F > 1$ (``hyper-falsity'') appear in nearly all vendors under the
appropriate protocol. These values are mathematically inconsistent
in Type-1 but are precisely representable in Type-2.

The rest of the paper is organised as follows. Section~2 reviews
Type-1 neutrosophic sets and existing extensions.
Section~3 introduces Type-$k$ Neutrosophic Sets (Definitions~1--3)
and states the strict hierarchy theorem. Section~4 defines the
recursive aggregation operator. Section~5 presents the empirical
evidence from LLM evaluations. Section~6 connects the construction
to existing neutrosophic and plithogenic frameworks.
Section~7 concludes.

% ====================================================================
\section{Background}

\subsection{Type-1 Neutrosophic Sets}

\begin{definition}[Type-1 Neutrosophic Set, \cite{smarandache1998neutrosophy}]
Let $U$ be a universe of discourse. A \emph{neutrosophic set}
$A^{(1)}$ on $U$ assigns to each element $x \in U$ a triplet
\[
  \mu^{(1)}(x) = (T, I, F), \quad T, I, F \in [0,1],
\]
where $T, I, F$ are independent: there is no constraint
$T + I + F = 1$.
The special case $T + I + F \leq 1$ recovers the
intuitionistic fuzzy set of Atanassov~\cite{atanassov1986ifs}.
\end{definition}

\subsection{Existing extensions}

Several extensions of Type-1 have been proposed.
\emph{Interval-valued neutrosophic sets}
\cite{wang2010svn} allow $T, I, F$ to be closed intervals
in $[0,1]$. \emph{Refined neutrosophic logic}
\cite{smarandache2013refined} splits each component into
sub-components: $T \to (T_1, \ldots, T_p)$,
$I \to (I_1, \ldots, I_r)$, $F \to (F_1, \ldots, F_s)$.
These sub-components remain scalars.
\emph{Neutrosophic overset / underset / offset}
\cite{smarandache2016overset} allows components outside
$[0,1]$: $T, I, F \in [-\Omega, \Omega^+]$ for bounds
$\Omega, \Omega^+ > 0$.
\emph{Plithogenic sets} \cite{smarandache2018plithogenic}
assign a contradiction degree between attribute values
and extend the aggregation accordingly.

None of these extensions assign a \emph{full neutrosophic
triplet} to each component, which is the defining feature
of Type-$k$ Neutrosophic Sets.

% ====================================================================
\section{Type-$k$ Neutrosophic Sets}

\subsection{Definitions}

\begin{definition}[Type-2 Neutrosophic Set]
\label{def:type2}
Let $U$ be a universe of discourse. A
\emph{Type-2 Neutrosophic Set} $A^{(2)}$ on $U$ assigns
to each $x \in U$ nine scalars organised as three nested
triplets:
\begin{align}
  \mu^{(2)}(x) &=
  \Bigl(
    \underbrace{(T_T,\, I_T,\, F_T)}_{\text{triplet of } T},\;
    \underbrace{(T_I,\, I_I,\, F_I)}_{\text{triplet of } I},\;
    \underbrace{(T_F,\, I_F,\, F_F)}_{\text{triplet of } F}
  \Bigr),
\end{align}
where all nine scalars are in $[0,1]$ independently.
The \emph{semantics} of each sub-triplet is:
\begin{itemize}
  \item $(T_T, I_T, F_T)$: to what degree is $T$ itself
    true / indeterminate / false.
  \item $(T_I, I_I, F_I)$: to what degree is $I$ itself
    true / indeterminate / false.
  \item $(T_F, I_F, F_F)$: to what degree is $F$ itself
    true / indeterminate / false.
\end{itemize}
\end{definition}

\begin{definition}[Type-3 Neutrosophic Set]
\label{def:type3}
A \emph{Type-3 Neutrosophic Set} $A^{(3)}$ on $U$ assigns
to each $x \in U$ twenty-seven scalars. Each of the nine
scalars of a Type-2 representation is itself replaced by
a neutrosophic triplet. For the $T$ component:
\begin{align}
  T \;\to\; &
  \Bigl(
    T_T(T_{TT}, I_{TT}, F_{TT}),\;
    I_T(T_{IT}, I_{IT}, F_{IT}),\;
    F_T(T_{FT}, I_{FT}, F_{FT})
  \Bigr),
\end{align}
and analogously for $I$ and $F$.
\end{definition}

\begin{definition}[Type-$k$ Neutrosophic Set]
\label{def:typek}
Let $k \geq 1$ be an integer.
A \emph{Type-$k$ Neutrosophic Set} $A^{(k)}$ on $U$ is
defined recursively:
\begin{itemize}
  \item \emph{Base case} ($k=1$): $A^{(1)}$ is a Type-1
    neutrosophic set (Definition~1).
  \item \emph{Recursive case} ($k \geq 2$): $A^{(k)}$ is
    obtained from $A^{(k-1)}$ by replacing each scalar
    component with a full neutrosophic triplet of scalars
    in $[0,1]$.
\end{itemize}
Each element $x \in U$ is characterised by $3^k$ scalars.
The set of Type-$k$ neutrosophic sets over $U$ is denoted
$\mathcal{NS}^{(k)}(U)$.
\end{definition}

\begin{remark}[Subset generalisation]
Definitions~\ref{def:type2}--\ref{def:typek} use scalar
components for finite depth $k$. A more general construction
allows $T$, $I$, $F$ to be neutrosophic \emph{sets} rather
than scalars, yielding an infinite-dimensional generalisation
that subsumes all finite Type-$k$ as projections. We leave
the formal development of this case to future work.
\end{remark}

\subsection{Canonical embedding}

\begin{definition}[Embedding $\varphi$]
The canonical embedding $\varphi: \mathcal{NS}^{(k-1)}(U)
\to \mathcal{NS}^{(k)}(U)$ maps each scalar $s$ in a
Type-$(k-1)$ element to the Type-$k$ triplet
$(s, 0, 1-s)$, i.e., a certain scalar is embedded as a
triplet with zero indeterminacy and complementary falsity.
\end{definition}

\subsection{Expressive hierarchy}

\begin{theorem}[Strict hierarchy]
\label{thm:hierarchy}
$\mathcal{NS}^{(1)}(U) \subsetneq \mathcal{NS}^{(2)}(U)
\subsetneq \cdots \subsetneq \mathcal{NS}^{(k)}(U)$
in expressive power.
\end{theorem}

\begin{proof}
Inclusion $\subseteq$: every Type-$(k-1)$ element $x$
embeds into Type-$k$ via $\varphi$ (Definition~4).

Strictness $\subsetneq$ (witness for $k=2$):
Consider the proposition ``This statement is both true and
false'' evaluated by a language model under an epistemic
protocol that allows values in $[-1, 2]$.
Suppose the model returns
$\mu^*_T = -0.20$, meaning ``the truth of this claim is
more false than true'' (anti-truth regime).
In Type-1, we require $T \in [0,1]$, so $\mu^*_T = -0.20$
is \emph{inadmissible}.
In Type-2, the same state is admissible as:
\[
  (T_T, I_T, F_T) = (0.00,\; 0.20,\; 0.80),
\]
meaning $T$ is 0\% true, 20\% indeterminate, and 80\% false
— a well-formed Type-2 triplet with all components in
$[0,1]$.
Empirically, we observe 65 such anti-truth responses
(2.7\% of evaluations; see Section~5).
The same argument generalises to the witness
$(T=1.7, I=1.8, F=1.6)$ (Table~\ref{tab:witnesses}),
which is inadmissible in any Type-1 extension with
$T,I,F \in [0,1]$ but fully representable in Type-2.
\end{proof}

% ====================================================================
\section{Recursive Aggregation Operator}

The standard SVNWA operator \cite{wang2010svn} at Type-1:
\[
  \text{SVNWA}(\mu_1,\ldots,\mu_n;\mathbf{w}) =
  \Bigl(1-\prod_{i=1}^n (1-T_i)^{w_i},\;
  \prod_{i=1}^n I_i^{w_i},\;
  \prod_{i=1}^n F_i^{w_i}\Bigr).
\]

\begin{definition}[Type-$k$ SVNWA]
The \emph{Type-$k$ SVNWA} applies the SVNWA operator
recursively to each sub-triplet independently at every
level. For Type-2 with weight vector $\mathbf{w}$:
\begin{align}
  \text{SVNWA}^{(2)}(\mu^{(2)}_1,\ldots,\mu^{(2)}_n;\mathbf{w}) &=
  \Bigl(
    \text{SVNWA}(T_1,\ldots,T_n;\mathbf{w}),\;
    \text{SVNWA}(I_1,\ldots,I_n;\mathbf{w}),\;
    \text{SVNWA}(F_1,\ldots,F_n;\mathbf{w})
  \Bigr),
\end{align}
where each argument of the outer tuple is itself the
SVNWA of the corresponding sub-triplets across elements
$1,\ldots,n$.
\end{definition}

\begin{proposition}[Reduction]
$\text{SVNWA}^{(k)}$ reduces to the standard SVNWA when
all Type-$k$ elements are images of Type-1 elements under
$\varphi$ (i.e., when all sub-triplets have $I=0$ and
$F=1-T$).
\end{proposition}

\begin{proof}
Under the embedding $\varphi$, each sub-triplet
$(s, 0, 1-s)$ has $I=0$, so $\prod I_i^{w_i} = 0$,
and $F = 1-T$ recovers the standard complement.
The SVNWA of the $T$-sub-triplets then equals
the standard SVNWA of the original $T$ values.
\end{proof}

% ====================================================================
\section{Empirical Evidence: LLMs as Implicit Type-2 Computers}

\subsection{Experimental design}

We reanalysed 2,419 evaluations from a multi-vendor
empirical study \cite{leyvavazquez2026plithogenic}.
Six foundation models were evaluated on five epistemic
phenomena (Paradox, Ignorance, Vagueness, Contradiction,
Contingency) under five protocols:
\begin{itemize}
  \item S1: classical single scalar $T \in [0,1]$.
  \item S4: Mason's \cite{mason2026tensors} three-component
    declared-loss frame $(what, why, severity)$.
  \item S4-N: per-attribute neutrosophic $(T,I,F)$ triplet.
  \item S4-O.A: overset protocol with $T, I, F \in [0,2]$.
  \item S4-O.C: offset protocol with $T, I, F \in [-1,2]$.
\end{itemize}
The six vendors were: Alibaba (Qwen-3-235B), Anthropic
(Claude Sonnet 4), DeepSeek Chat, Meta (Llama-4-Maverick),
Mistral Medium-3.1, and OpenAI (GPT-4o).

\subsection{Type-2 identification criterion}

A response $\mu = (T, I, F)$ \emph{requires Type-2
representation} if and only if at least one component
falls outside $[0,1]$:
$T < 0$, $T > 1$, $I > 1$, or $F > 1$.
Such responses are inadmissible in any Type-1 neutrosophic
set (including the overset/offset extensions of
\cite{smarandache2016overset} when the extended range
is \emph{not} assumed) but are representable in Type-2
via the explicit mapping defined in the proof of
Theorem~\ref{thm:hierarchy}.

\subsection{Results}

\begin{table}[h!]
\centering
\caption{Type-2 requirement by protocol ($n$ = evaluations per protocol).}
\label{tab:protocol}
\begin{tabular}{lrrr}
\toprule
Protocol & $n$ & Type-2 required & \% \\
\midrule
S4-O.A  & 300 & 292 & \textbf{97.3} \\
S4-O.C  & 889 & 305 & \textbf{34.3} \\
S4-N    & 300 & 0   & 0.0 \\
S4      & 300 & 0   & 0.0 \\
S1      & 300 & 0   & 0.0 \\
S1-taut &  90 & 0   & 0.0 \\
\midrule
\textbf{Total} & \textbf{2,419} & \textbf{597} & \textbf{24.7} \\
\bottomrule
\end{tabular}
\end{table}

\begin{table}[h!]
\centering
\caption{Type-2 requirement by vendor (S4-O protocols only).}
\label{tab:vendor}
\begin{tabular}{lrrr}
\toprule
Vendor & $n$ & Type-2 required & \% \\
\midrule
Mistral (Medium-3.1)         & 435 & 130 & 29.9 \\
OpenAI (GPT-4o)              & 377 & 100 & 26.5 \\
Anthropic (Claude Sonnet 4)  & 408 & 108 & 26.5 \\
Meta (Llama-4-Maverick)      & 355 &  81 & 22.8 \\
DeepSeek Chat                & 431 &  97 & 22.5 \\
Alibaba (Qwen-3-235B)        & 413 &  81 & 19.6 \\
\bottomrule
\end{tabular}
\end{table}

\begin{table}[h!]
\centering
\caption{Type-2 subtypes across all evaluations.}
\label{tab:subtypes}
\begin{tabular}{lrr}
\toprule
Type-2 subtype & $n$ & \% of total \\
\midrule
$I > 1$ (hyper-indeterminacy)  & 507 & 21.0 \\
$T > 1$ (hyper-truth)          & 129 &  5.3 \\
$F > 1$ (hyper-falsity)        &  91 &  3.8 \\
$T < 0$ (anti-truth)           &  65 &  2.7 \\
\bottomrule
\end{tabular}
\end{table}

\begin{table}[h!]
\centering
\caption{Five most extreme Type-2 witnesses.}
\label{tab:witnesses}
\begin{tabular}{lllccc}
\toprule
Vendor & Protocol & Phenomenon & $T$ & $I$ & $F$ \\
\midrule
Mistral & S4-O.A & Contradiction (Ethical) & 1.7 & 1.8 & 1.6 \\
Mistral & S4-O.A & Contradiction (Ethical) & 1.7 & 1.8 & 1.6 \\
OpenAI  & S4-O.A & Paradox (Logical)       & 0.0 & 2.0 & 2.0 \\
OpenAI  & S4-O.A & Paradox (Logical)       & 0.0 & 2.0 & 2.0 \\
OpenAI  & S4-O.A & Paradox (Logical)       & 0.0 & 2.0 & 2.0 \\
\bottomrule
\end{tabular}
\end{table}

\subsection{Interpretation}

Three observations are central:

\textbf{1. Zero vs.\ high dichotomy by protocol.}
Standard protocols (S1, S4, S4-N) produce \emph{zero}
Type-2 responses, not because the epistemic phenomena
do not require Type-2, but because the protocol's
output format enforces $T, I, F \in [0,1]$.
Extended protocols (S4-O.A, S4-O.C) release this
constraint, and nearly all models immediately occupy
the Type-2 region (97.3\% under S4-O.A).
This demonstrates that the Type-1 restriction is a
\emph{protocol artefact}, not a property of the
epistemic phenomena being evaluated.

\textbf{2. Cross-vendor consistency.}
The range 19.6\%--29.9\% across six independent
foundation models (Table~\ref{tab:vendor}) rules out
the possibility that Type-2 behaviour is an idiosyncrasy
of a single vendor's training. All six models, trained
by different organisations on different data with
different objectives, exhibit the same qualitative
behaviour.

\textbf{3. Type-2 mapping of extreme witnesses.}
The witness $(T=1.7, I=1.8, F=1.6)$ from Mistral
is inadmissible as a Type-1 value. Its Type-2
representation is:
\begin{align}
  T \;\to\; (T_T=1.0,\; I_T=0.7,\; F_T=0.0), \nonumber\\
  I \;\to\; (T_I=1.0,\; I_I=0.8,\; F_I=0.0), \nonumber\\
  F \;\to\; (T_F=1.0,\; I_F=0.6,\; F_F=0.0). \nonumber
\end{align}
Under this mapping, the model is expressing:
``truth is maximally true but 70\% additionally
indeterminate; indeterminacy is maximally true but
80\% additionally indeterminate.'' This is a coherent
Type-2 epistemic state that Type-1 cannot represent.

% ====================================================================
\section{Connections to Existing Neutrosophic Frameworks}

\subsection{Refined neutrosophic logic}

Refined neutrosophic logic \cite{smarandache2013refined}
increases the \emph{cardinality} at a single level
($T \to T_1, \ldots, T_p$).
Type-$k$ Neutrosophic Sets increase the \emph{depth}
of the representation.
The two constructions are orthogonal and composable:
a Type-2 refined set would split each sub-triplet
component into sub-components at the same level.

\subsection{Plithogenic tensors}

Leyva-V\'{a}zquez and Smarandache \cite{leyvavazquez2026plithogenic}
assign independent $(T, I, F)$ triplets to each
\emph{attribute} of the object being evaluated.
This is a Type-2 structure along the attribute
dimension: the overall epistemic state of an attribute
is itself characterised by a triplet.
Type-$k$ Neutrosophic Sets generalise this to arbitrary
recursion depth.

\subsection{SVN Tensors}

Single-valued neutrosophic tensors
\cite{leyvavazquez2026svntensors} are Type-1 structures
arranged along multiple tensor modes.
Every SVN tensor is a special case of a Type-1 element
under the framework of this paper.
Type-$k$ extends SVN tensors to depth $k$ by replacing
each scalar component with a sub-triplet.

\subsection{The Absorption Problem}

Prior work \cite{leyvavazquez2026plithogenic} identified
the \emph{Absorption Problem}: under Type-1 protocols,
indeterminacy $I$ systematically absorbs the entire
probability mass ($I \to 1$) when the phenomenon is
genuinely uncertain.
The Type-2 framework reframes this as a depth-collapse:
the Absorption Problem is forced Type-2 $\to$ Type-1
projection in which the sub-structure of $I$
(namely, whether $I$ is high due to genuine ignorance
vs.\ due to source conflict) is lost.
Type-2 representations distinguish $(T_I, I_I, F_I)
= (1, 0, 0)$ (certain indeterminacy) from
$(T_I, I_I, F_I) = (0, 0, 1)$
(the apparent indeterminacy is itself false, i.e.,
the state is actually a contradiction).
These two states are identical under Type-1 collapse.

% ====================================================================
\section{Conclusion}

We have introduced Type-$k$ Neutrosophic Sets as a
formal recursive extension of neutrosophic logic,
proved a strict expressive hierarchy, and demonstrated
empirically that state-of-the-art LLMs already compute
Type-2 values when the evaluation protocol allows it.

Three directions for future work are immediate:
\begin{enumerate}
  \item \textbf{Type-3 for meta-uncertainty.}
    A Type-3 evaluation would allow the evaluator to
    characterise its own indeterminacy: not only
    ``how uncertain am I?'' (captured in $I$) but
    ``how certain am I about my uncertainty?''
    (captured in $I_I$ at Type-3).
  \item \textbf{Subset generalisation.}
    Replacing scalar components with neutrosophic
    sets (rather than finite recursive triplets)
    yields an infinite-dimensional framework that
    subsumes all finite Type-$k$ as projections.
  \item \textbf{Type-$k$ distance and clustering.}
    A Type-$k$ distance metric would enable clustering
    of LLMs by their characteristic depth of epistemic
    expression, potentially distinguishing models that
    compute genuine Type-2 uncertainty from those
    that produce numerical artefacts.
\end{enumerate}

\section*{Acknowledgements}
The authors thank the six foundation model providers
for their publicly accessible APIs.

% ====================================================================
\begin{thebibliography}{99}

\bibitem{smarandache1998neutrosophy}
Smarandache, F. (1998).
\textit{Neutrosophy: Neutrosophic Probability, Set, and Logic}.
American Research Press.

\bibitem{smarandache2013refined}
Smarandache, F. (2013).
n-Valued refined neutrosophic logic and its applications
to physics.
\textit{Progress in Physics}, 4, 143--146.

\bibitem{smarandache2016overset}
Smarandache, F. (2016).
Neutrosophic overset, neutrosophic underset, and
neutrosophic offset.
\textit{Pons Editions}.

\bibitem{smarandache2018plithogenic}
Smarandache, F. (2018).
Plithogeny, plithogenic set, logic, probability, and
statistics.
\textit{Pons Publishing House}.

\bibitem{wang2010svn}
Wang, H., Smarandache, F., Zhang, Y.\,Q., \& Sunderraman, R. (2010).
Single valued neutrosophic sets.
\textit{Multispace \& Multistructure}, 4, 410--413.

\bibitem{atanassov1986ifs}
Atanassov, K.\,T. (1986).
Intuitionistic fuzzy sets.
\textit{Fuzzy Sets and Systems}, 20(1), 87--96.

\bibitem{mendel2002type2}
Mendel, J.\,M., \& John, R.\,I. (2002).
Type-2 fuzzy sets made simple.
\textit{IEEE Transactions on Fuzzy Systems}, 10(2), 117--127.

\bibitem{mason2026tensors}
Mason, T. (2026).
From scalars to tensors: Declared losses recover
epistemic distinctions that neutrosophic scalars
cannot express.
arXiv:2604.09602v1 [cs.AI].

\bibitem{leyvavazquez2026plithogenic}
Leyva-V\'{a}zquez, M., \& Smarandache, F. (2026).
Plithogenic tensor neutrosophic logic: A formal
foundation for declared-loss evaluation in large
language models.
\textit{Manuscript submitted}.

\bibitem{leyvavazquez2026svntensors}
Leyva-V\'{a}zquez, M., \& Smarandache, F. (2026).
Single-valued neutrosophic tensors.
\textit{arXiv preprint} (in preparation).

\bibitem{ye2014svnwa}
Ye, J. (2014).
A multicriteria decision-making method using
aggregation operators for simplified neutrosophic sets.
\textit{Journal of Intelligent \& Fuzzy Systems},
26(5), 2459--2466.

\end{thebibliography}

\end{document}
